\PassOptionsToPackage{colorlinks=true,allcolors=black}{hyperref}
% REQ: The above line, \PassOptionsToPackage{} must be very first line.

\documentclass[11pt]{article}

% Encoding and font setup for cross-platform reproducibility
\usepackage[T1]{fontenc}
\usepackage{lmodern}

% Improve readability for dense theoretical material
\usepackage{setspace}
\onehalfspacing

% Mathematical notation and table formatting for formal semantics
\usepackage{amsmath, amssymb}
\usepackage{array}
\usepackage[table]{xcolor}
\usepackage{colortbl}
\usepackage{booktabs}
\usepackage{graphicx}

% Theorem environments
\usepackage{amsthm}

% =============================================================================
% Numbering convention:
% - All theorem-like statements share one counter, numbered by section.
% - Prose cites statements by stable names and labels, not by number.
% - External documents cite by stable tag (e.g., se100.def.Substrate).
%   Numbers move when text is edited; names do not.
% =============================================================================
\theoremstyle{plain}
\newtheorem{theorem}{Theorem}[section]
\newtheorem{lemma}[theorem]{Lemma}
\newtheorem{proposition}[theorem]{Proposition}
\newtheorem{constraint}[theorem]{Constraint}
\newtheorem{claim}[theorem]{Claim}
\newtheorem{assumption}[theorem]{Assumption}

\theoremstyle{definition}
\newtheorem{definition}[theorem]{Definition}
\newtheorem{example}[theorem]{Example}

\theoremstyle{remark}
\newtheorem{remark}[theorem]{Remark}
\newtheorem{note}[theorem]{Note}

% Unnumbered versions for front-matter previews/callouts
\newtheorem*{constraint*}{Constraint}
\newtheorem*{definition*}{Definition}
\newtheorem*{example*}{Example}
\newtheorem*{remark*}{Remark}
\newtheorem*{note*}{Note}

% Bibliography management and hyperlink support
\usepackage{url}
\usepackage{natbib}
\bibpunct{(}{)}{;}{a}{}{ }

% Core hyperlink and bookmark infrastructure
\usepackage{hyperref}
\usepackage{bookmark}

% Support multiple author affiliations
\usepackage{authblk}

% Improve typographic quality and line breaking
\usepackage{microtype}

% Callout boxes for examples and emphasis
\usepackage[most]{tcolorbox}

% Compact list formatting for dense conceptual content
\usepackage{enumitem}
\setlist[itemize]{itemsep=0.2em, topsep=0.2em, parsep=0em, partopsep=0em}
\setlist[enumerate]{itemsep=0.2em, topsep=0.2em, parsep=0em, partopsep=0em}

% Paragraph spacing prioritizes readability over traditional indentation
\setlength{\parskip}{0.75em}
\setlength{\parindent}{0em}

% Keywords macro for structured abstract metadata
\providecommand{\keywords}[1]{\vspace{0.5em}\noindent\textbf{Keywords: }\small #1}

% =============================================================================
% Stable named-reference helpers.
% Use labels such as se100.def.Substrate and cite by stable name:
% \DefRef{se100.def.Substrate}{Substrate}
% =============================================================================
\newcommand{\DefRef}[2]{\hyperref[#1]{Definition (#2)}}
\newcommand{\AssumpRef}[2]{\hyperref[#1]{Assumption (#2)}}
\newcommand{\ConstRef}[2]{\hyperref[#1]{Constraint (#2)}}
\newcommand{\ExRef}[2]{\hyperref[#1]{Example (#2)}}
\newcommand{\RemRef}[2]{\hyperref[#1]{Remark (#2)}}
\newcommand{\NoteRef}[2]{\hyperref[#1]{Note (#2)}}
\newcommand{\SecRef}[1]{Section~\ref{#1}}

% =============================================================================
% Formal notation
% =============================================================================
\newcommand{\Substrate}{\mathcal{S}}
\newcommand{\SubstrateRef}{\mathcal{S}_{\mathrm{ref}}}
\newcommand{\Frameworks}{\mathbb{F}}
\newcommand{\Framework}{F}
\newcommand{\FrameworkOne}{F_1}
\newcommand{\FrameworkTwo}{F_2}
\newcommand{\entails}{\vdash}
\newcommand{\notentails}{\nvdash}
\newcommand{\Asserts}{\mathsf{Asserts}}
\newcommand{\pred}[1]{\mathit{#1}}

% Reusable callout box
\newcommand{\FigureCallout}[2]{%
  \begin{tcolorbox}[
    colback=gray!5,
    colframe=black!40,
    title={#1},
    fonttitle=\bfseries,
    arc=3pt,
    boxrule=0.5pt,
    width=\linewidth,
    enhanced,
    breakable
  ]
  #2
  \end{tcolorbox}
}

% =============================================================================
% PREFERRED CITATION SPECIFICATION
%
% @techreport{case2026neutral,
%   author      = {Case, Denise M.},
%   title       = {Neutral Substrates: A Design Constraint for Shared Records Under Persistent Interpretive Disagreement},
%   institution = {Preprint},
%   year        = {2026},
%   url         = {}
% }
% =============================================================================

% ============================================================
\title{
Neutral Substrates:\\
A Design Constraint for Shared Records\\
Under Persistent Interpretive Disagreement
}
% ============================================================

\author{Denise M. Case}
\affil{\small Department of Computer Science and Information Systems\\
Northwest Missouri State University, Maryville, MO, USA\\}
\date{}

\begin{document}
\maketitle
\vspace{-2em}

\begin{abstract}
  Shared accountability records are often used
  by parties who do not, and may never,
  agree about causation, responsibility, or normative interpretation.
  In such settings, a shared record cannot remain neutral
  by omitting contested information,
  because accountability requires preserving the claims parties made,
  together with the sources and provenance on which those claims rely.
  Nor can it remain neutral
  by asserting one contested interpretation as the shared base.

  This paper defines a neutral substrate as a shared representational layer
  that provides stable reference
  while making no object-level substrate-layer commitments
  to causal or normative propositions.
  The central design constraint is that,
  when causal and normative propositions are
  contestable across admissible frameworks and
  the substrate's referential commitments
  are common ground among those frameworks,
  a substrate satisfies neutrality by design
  if and only if its foundational layer is restricted to
  those referential commitments and attribution propositions
  whose attributional basis is fixed by those referential commitments.
  Causal and normative content may still be represented,
  but not as object-level foundational-layer commitments:
  it may appear there only as the content of attributed assertions with provenance,
  made by some identified framework, source, agent, institution, record, or document.

  The representational machinery used here is standard:
  reification, attribution, and provenance.
  The contribution is the constraint:
  a precise, checkable condition on the foundational layer of a shared record,
  stated together with the assumptions it depends on and the boundary condition
  under which it fails.
  A neutral substrate says enough to preserve accountability,
  but it does not turn one party's interpretation
  into an object-level substrate-layer commitment.
  Neutrality is unavailable at that layer when the referential regime is
  contested among the frameworks in play.
\end{abstract}


\keywords{
  Formal ontology;
  neutral substrates;
  accountable records;
  referential regimes;
  interpretive non-commitment;
  extension stability;
  persistent interpretive disagreement;
  reification;
  provenance
}

\begin{note*}[Citation Convention]
  Definitions, assumptions, constraints, and examples in this paper carry stable
  names (for example, \texttt{se100.def.Substrate}).
  This paper cites statements by stable name rather than by number,
  so citations remain valid even if numbering changes.
\end{note*}

% =============================================================================
% Preview of Main Constraint
% =============================================================================

\bigskip
\noindent\textbf{Preview of Main Constraint.}
\emph{
  Let $\Substrate$ be a substrate
  intended to remain usable by every admissible interpretive framework,
  including admissible frameworks not known when
  $\Substrate$ is designed.
  Suppose causal and normative propositions are
  contestable across admissible frameworks,
  and suppose the referential commitments of $\Substrate$
  are common ground among those frameworks.
  Then $\Substrate$ is neutral by design
  if and only if its foundational layer is restricted to
  the substrate's referential commitments
  and attribution propositions
  whose attributional basis is fixed by those referential commitments.
  In particular, $\Substrate$ makes no substrate-layer commitment
  to any object-level causal or normative proposition.
}

\medskip
\noindent
\emph{Informally.} A neutral substrate may preserve causal and normative claims,
but it must preserve them as attributed, provenance-bearing assertions,
not as object-level substrate-layer commitments.

% =============================================================================
\section{Introduction}
\label{sec:introduction}
% =============================================================================

Accountability records are rarely interpreted by only one party.
A record may be used by institutions, auditors, affected individuals,
regulators, courts, researchers, vendors, advocates, or future maintainers.
These parties may share enough reference to identify
the person, event, decision, instrument, place, file, policy, rule,
or institutional artifact at issue.
At the same time, they may disagree, possibly permanently,
about what caused what, who was responsible,
whether a decision was justified, which rule applied,
or which normative standard should control.

This creates a structural design problem.
A shared record must remain usable by multiple interpretive frameworks,
including frameworks that permanently disagree.
If the shared layer makes a substrate-layer commitment
to a contested causal or normative proposition,
the record becomes structurally aligned with one framework against another.
It may still be useful to the favored framework, but it no
longer functions as a neutral substrate for shared accountability.

The purpose of this paper is to define the substrate-layer constraint
needed to avoid that failure.

The central claim is:

\begin{quote}
  A neutral substrate may preserve causal or normative claims,
  but it may not adopt them as object-level substrate-layer commitments.
  It must represent them as attributed, provenance-bearing assertions
  about the referents fixed by the substrate.
\end{quote}

The mechanism of neutrality is the distinction between
asserting a proposition and representing that some identified source $x$
asserted it.
This mechanism avoids two failures.
It does not suppress contested content,
because the attributed claim remains in the record.
It does not adopt contested content,
because the substrate commits to the attribution proposition
$\Asserts(x,\varphi)$,
not to the asserted proposition $\varphi$.

The mechanism is not new (\SecRef{sec:related-work}).
The contribution is the constraint:
a biconditional connecting neutrality by design
to a checkable restriction on the substrate's foundational layer,
together with the assumptions and boundary condition on which the constraint depends.

The paper proceeds as follows.
\SecRef{sec:related-work} credits the foundations and states the contribution.
\SecRef{sec:problem-setting} states the problem.
\SecRef{sec:definitions} gives the core definitions.
\SecRef{sec:neutrality} states and justifies the neutrality constraint.
\SecRef{sec:mechanism} explains the reification mechanism.
\SecRef{sec:example} gives a worked example.
\SecRef{sec:boundary} states the boundary condition under which neutrality is unavailable.
\SecRef{sec:non-goals} identifies non-goals.
\SecRef{sec:discussion} discusses the practical significance and limits of the constraint.
\SecRef{sec:conclusion} concludes.

% =============================================================================
\section{Related Work and Contribution}
\label{sec:related-work}
% =============================================================================

The representational mechanism this paper relies on is not novel.
Prior work developed machinery for
recording that a source asserts a proposition
rather than asserting the proposition itself.
This is standard in knowledge representation,
logical AI,
argumentation theory,
and applied ontology.

In knowledge representation on the web,
RDF reification \citep{hayes2004rdf}
and named graphs \citep{carroll2005named}
provide standard machinery often used for
statements about statements.
The W3C provenance ontology
\citep{lebo2013prov} standardizes attribution of content to
agents, activities, and sources.
The form $\Asserts(x,\varphi)$ used in this paper,
with $x$ identifying the source to which the assertion is attributed,
follows that representational pattern;
nothing representational is added.

In logical AI, \citet{mccarthy1993} treats truth as relative to a context via
$\mathit{ist}(c,p)$.
Framework-relative derivability
($\Substrate \cup \Framework \entails p$)
uses the same general idea in a narrower way.
Purpose is the difference:
context logic supports reasoning across contexts,
whereas the substrate defined here deliberately abstains from
cross-framework commitments at its foundational layer.
Assumption-based truth maintenance \citep{dekleer1986atms}
likewise maintains conclusions relative to assumption sets
and keeps mutually inconsistent worlds simultaneously available.
Extension stability, defined below,
is a record-level analogue of that discipline:
an ATMS is a reasoning architecture,
whereas this constraint is a design condition on records.

In argumentation theory, \citet{dung1995} showed how to compute acceptability
among conflicting arguments.
The aim here is complementary and prior:
argumentation frameworks adjudicate among conflicting claims,
whereas a neutral substrate refuses adjudication at the base layer,
preserving the conflicting claims with attribution.
A neutral substrate is the kind of shared input over
which a Dung-style evaluation could later be run by a framework,
not by the substrate.

Two older distinctions underwrite the whole construction.
Speech act theory distinguishes asserting a proposition
from reporting that someone asserted it \citep{searle1969}.
The analysis of ontological commitment \citep{guarino1998}
makes precise the idea that a representation commits its
users to a view of the domain.
The substrate-layer commitment defined below is
a record-level application of that idea.

Within applied ontology,
the descriptions-and-situations pattern of DOLCE+DnS Ultralite
represents descriptions, situations, and interpretive stances
as first-class modeling elements distinct
from the states of affairs they describe
\citep{gangemi2003descriptions}.
The present constraint is compatible with that pattern and
can be read as a condition on which layer a commitment occupies:
reference and permitted attribution at the substrate,
causal and normative interpretation in framework-dependent descriptions.
What this paper adds is not a new reification mechanism,
but a design-time restriction on which commitments may
belong to the shared foundational layer
when neutrality under persistent interpretive disagreement is required.

This paper is also adjacent to Floridi's method of levels of abstraction
\citep{floridi2008loa}.
The analysis here is stated as a representational-layer constraint
rather than as a level-of-abstraction claim:
the issue is not only which abstraction an interpreter selects,
but which commitments the shared record itself adopts
at its foundational layer.

\textbf{The contribution.}
Prior literatures provide the ingredients but do not state
a checkable condition on the foundational layer of a shared accountability record.
This paper assembles those ingredients into a single biconditional constraint
on the substrate layer.

Reification and provenance provide the mechanism.
Context logic and truth maintenance provide the relativized-derivability reading.
Argumentation theory provides the downstream consumer.
Speech act theory and ontological commitment provide the underlying distinctions.

The constraint combines an admissibility gate for frameworks,
a contestability assumption for causal and normative content,
a common-ground assumption for reference,
and a pairwise stability condition.
It also states the admissibility and grounding conditions that make it usable,
including the exact boundary at which it fails.

Each ingredient is known.
The addition is the constraint and the discipline of saying precisely
when it does not apply.

% =============================================================================
\section{Problem Setting}
\label{sec:problem-setting}
% =============================================================================

A shared record can fail neutrality in two opposed ways.

\begin{enumerate}
  \item It can omit accountability-relevant claims entirely.
        This avoids adopting those claims as object-level substrate-layer commitments,
        but it also suppresses the information needed to understand
        who claimed what, under what authority, and on what basis.
        A record that suppresses contested claims cannot support accountability
        for those claims.

  \item It can make a substrate-layer commitment
        to a contested causal or normative proposition.
        This preserves the content of the claim,
        but strips away the attribution that makes the claim contestable.
        Such a record may appear complete, but is structurally non-neutral.
\end{enumerate}

A neutral substrate must avoid both failures.
It must preserve contested claims without adopting them
as object-level substrate-layer commitments.

This paper uses a role-based taxonomy of \emph{substrate-carried content}:

\begin{enumerate}
  \item \textbf{Referential content}: entities, occurrences, institutional
        artifacts, identifiers, timestamps, provenance, and the relations needed
        to fix what is being referred to.
  \item \textbf{Attribution content}: propositions of the form
        $\Asserts(x,\varphi)$, meaning that some
        framework, source, agent, institution, record, or document $x$
        asserts a claim $\varphi$ about the referents.
  \item \textbf{Object-level interpretive content}: the asserted propositions
        $\varphi$ themselves about the referents, rather than attribution
        propositions about those assertions.
        In this paper, the neutrality constraint applies to object-level causal or normative propositions.
\end{enumerate}

This taxonomy distinguishes representational roles,
not mutually exclusive kinds of propositions.
When the substrate commits to content of any of these kinds,
that commitment is a substrate-layer commitment in the sense of
\DefRef{se100.def.SubstrateCommitment}{Substrate-Layer Commitment}.
The substrate may entail a referential proposition,
an attribution proposition,
or an object-level interpretive proposition.
The three kinds differ in \emph{what} is committed,
not in \emph{whether} the substrate commits.

Neutrality permits substrate-layer commitments to referential content
and attribution content.
It prohibits substrate-layer commitments,
made on the substrate's own account,
to contested object-level causal or normative content.
Such content must enter the record as attribution content,
not as an object-level causal or normative commitment of the substrate.

For neutrality by design, causal and normative content must not enter the
foundational layer as object-level substrate-layer commitments.
Such content may be preserved in the record only as attributed assertions
whose attributional basis, including provenance, is fixed by the substrate's
referential commitments.

% =============================================================================
\section{Definitions}
\label{sec:definitions}
% =============================================================================

% =============================================================================
\subsection*{Substrates and Commitments}
% =============================================================================

\begin{definition}[Substrate]
  \label{se100.def.Substrate}
  A \emph{substrate} $\Substrate$ is a shared representational base providing
  stable reference for entities, occurrences, and institutional artifacts
  across a class of admissible interpretive frameworks $\Frameworks$.
  Stable reference consists of individuation, co-reference, and persistence.
\end{definition}

A substrate is not the whole record system.
It is the foundational layer on which further interpretive claims may be layered,
including causal or normative interpretations,
evidentiary assessments, explanatory accounts, and institutional determinations.
Its purpose is to make shared reference possible without
forcing agreement about every interpretation of what is referenced.
A substrate may identify that there is an event, entity, decision,
document, instrument, rule, policy, observation, or institutional artifact, and
it may preserve identifiers, timestamps, provenance, and persistence relations
among them.
Whether some event caused another,
whether some decision was justified, or
whether some rule ought to apply
are not commitments of the substrate
just because the substrate identifies
the relevant event, decision, rule,
or institutional artifact.

% ==================== Substrate-Layer Commitment ============

\begin{definition}[Substrate-Layer Commitment]
  \label{se100.def.SubstrateCommitment}
  A substrate $\Substrate$ \emph{commits} to a proposition $p$ when
  $\Substrate \entails p$;
  that is, $p$ is asserted by the substrate independently of any interpretive
  framework.
\end{definition}

% =============================================================================
\subsection*{Proposition Roles}
% =============================================================================

\begin{note}[Causal and Normative Content as Primitive Classifications]
  \label{se100.note.CausalNormative}
  This paper treats \emph{causal proposition} and \emph{normative proposition}
  as primitive content classifications,
  not as terms reduced here to a formal decision procedure.
  As a working guide,
  a proposition is causal if it asserts that one event,
  state, action, condition, or process brought about, produced, prevented, or
  contributed to another.
  A proposition is normative if it asserts that something was justified,
  permitted, required, prohibited, correct, incorrect, compliant, or in
  violation of a rule, standard, policy, or norm.

  The classification is not by surface vocabulary.
  An innocuous-looking field can encode a causal or normative commitment,
  and causal or normative words can appear inside an attributed assertion
  without the substrate committing to the asserted causal or normative
  proposition.
  The constraint applies once the relevant content classification is fixed
  for the accountability context.
\end{note}



% ============  Attribution Propositions ============

\begin{definition}[Attribution Proposition]
  \label{se100.def.AttributionProposition}
  An \emph{attribution proposition} is a proposition of the form
  $\Asserts(x,\varphi)$,
  meaning that some framework, source, agent, institution, record, or document $x$
  asserts proposition $\varphi$.

  A substrate-layer commitment to $\Asserts(x,\varphi)$ is a
  commitment to the attribution,
  not to the asserted proposition $\varphi$.
\end{definition}



% ==================== Object-Level Interpretive Propositions ============

\begin{definition}[Object-Level Interpretive Proposition]
  \label{se100.def.ObjectLevelInterpretiveProposition}
  An \emph{object-level interpretive proposition} is an asserted proposition
  $\varphi$ itself,
  about the referents fixed by the substrate,
  rather than an attribution proposition of the form
  $\Asserts(x,\varphi)$.

  An \emph{object-level causal or normative proposition} is
  an object-level interpretive proposition whose content is causal or normative
  (\NoteRef{se100.note.CausalNormative}{Causal and Normative Content as
    Primitive Classifications}).
\end{definition}

The core distinction is that the substrate may commit to
$\Asserts(x,\varphi)$
without committing to $\varphi$.
Neutrality does not require suppressing causal or normative information.
It requires keeping causal and normative information at the level
of attribution propositions rather than substrate-layer commitments
to the asserted propositions.

\begin{definition}[Object-Level Causal or Normative Commitment]
  \label{se100.def.ObjectLevelCausalNormativeCommitment}
  An \emph{object-level causal or normative commitment} is a substrate-layer
  commitment to an object-level causal or normative proposition,
  rather than to an attribution proposition about that proposition.
\end{definition}


% =============================================================================
\subsection*{Reference}
% =============================================================================

\begin{definition}[Referential Regime]
  \label{se100.def.ReferentialRegime}
  A \emph{referential regime} is the triple of
  \begin{enumerate}
    \item individuation conditions,
    \item co-reference conditions, and
    \item persistence conditions
  \end{enumerate}
  by which a substrate fixes and tracks entities, occurrences, and institutional
  artifacts.
\end{definition}

Individuation conditions determine when something counts as one thing rather
than another.
Co-reference conditions determine when two references refer to the same thing.
Persistence conditions determine when something remains the same thing across
time, transformation, revision, or institutional change.
The referential regime is what allows a substrate to support shared use:
if parties cannot share enough reference to identify what is being interpreted,
they cannot reliably disagree about causes, norms, explanations, or obligations
concerning that thing.

\begin{definition}[Referential Commitments]
  \label{se100.def.ReferentialCommitments}
  Let $\SubstrateRef$ denote the \emph{referential commitments} of $\Substrate$.
  These are the substrate-layer commitments fixed by its referential regime,
  including identifiers;
  the typing of entities, occurrences, and institutional artifacts;
  timestamps;
  provenance;
  and referential relations among them.
\end{definition}

% =============================================================================
\subsection*{Frameworks}
% =============================================================================

\begin{definition}[Admissible Framework]
  \label{se100.def.AdmissibleFramework}
  A framework $\Framework$ is \emph{admissible} if it satisfies three conditions:
  \begin{enumerate}
    \item \textbf{Internal consistency}: $\Framework \notentails \bot$.
    \item \textbf{Evidentiary grounding}: where $\Framework$ makes empirical,
          causal, normative, or other interpretive claims, it identifies the evidence,
          source, method, measurement, observation, record, rule, standard, or
          document on which those claims rely.
    \item \textbf{Documented interpretive function}: $\Framework$ is presented as
          an interpretive function with a
          named source, documented scope, and citable basis.
  \end{enumerate}
\end{definition}

Admissibility is not agreement.
Two admissible frameworks may disagree about
causation, responsibility, justification, applicability, or normative significance.
Admissibility does not determine which framework should prevail.
It specifies the minimum conditions under which a framework can enter
the shared record as an identifiable, scoped, and citable interpretive position.

The admissibility gate is not neutral, and this paper does not claim that.
Requiring evidentiary grounding and a documented interpretive function
(Conditions~2 and~3) embeds design choices about what counts as
a framework that can enter the shared record.

These conditions do not adjudicate the framework's claims.
They make the framework available for representation in the record.
Any further claim about whether a framework's basis is sufficient,
whether its interpretation should govern,
or whether its use is appropriate in a particular dispute
must be represented as an attributed claim with provenance,
not as a substrate-layer commitment.

\begin{note}[The Framework Class $\Frameworks$]
  \label{se100.note.FrameworkClass}
  Throughout, $\Frameworks$ denotes the
  class of \emph{all} admissible interpretive frameworks.
  Admissibility is an intrinsic property
  (\DefRef{se100.def.AdmissibleFramework}{Admissible Framework}),
  so membership in $\Frameworks$ does not depend on designer knowledge.
  The class is open and not enumerable at design time:
  new admissible interpretive frameworks can be authored at any
  point in a record's lifetime.
  ``Extension of $\Frameworks$'' below refers to
  such frameworks becoming known and entering consideration,
  not to change in the intrinsic class.
  This sense of extension is distinct from the
  substrate-with-framework extension $\Substrate \cup \Framework$ named in
  \DefRef{se100.def.ExtensionStability}{Extension Stability}.
\end{note}

\begin{definition}[Permitted Attribution Proposition]
  \label{se100.def.PermittedAttributionProposition}
  An attribution proposition $\Asserts(x,\varphi)$ is
  \emph{permitted at the foundational layer} if the attributional basis for
  $x$'s assertion of $\varphi$ is fixed by the substrate's referential
  commitments $\SubstrateRef$.

  The attributional basis includes the source, assertion occurrence,
  provenance, and content reference needed to identify what was asserted,
  by whom, and under what record basis.

  A permitted attribution proposition commits the substrate
  to the attribution by $x$, not to the asserted proposition $\varphi$.
\end{definition}

\begin{definition}[Framework-Variant Proposition]
  \label{se100.def.FrameworkVariant}
  A proposition $p$ is \emph{framework-variant}
  with respect to substrate $\Substrate$ and framework class $\Frameworks$
  if there exist admissible frameworks $\FrameworkOne, \FrameworkTwo \in \Frameworks$
  such that
  \[
    \Substrate \cup \FrameworkOne \entails p
    \quad\text{and}\quad
    \Substrate \cup \FrameworkTwo \entails \neg p
  \]
\end{definition}

Framework variance captures disagreement that arises
not from different referential bases,
but from different admissible interpretations
layered on top of the same referential base.
A proposition accepted under one framework
and rejected under another is
not suitable as an unqualified substrate-layer commitment.

\begin{definition}[Framework-Invariant Proposition]
  \label{se100.def.FrameworkInvariant}
  A proposition $p$ is \emph{framework-invariant} with respect to substrate
  $\Substrate$ and framework class $\Frameworks$
  if, for every admissible framework $\Framework \in \Frameworks$,
  \[
    \Substrate \cup \Framework \cup \{p\} \notentails \bot .
  \]
  Equivalently,
  $p$ can be added to the substrate and
  remain compatible with every admissible framework:
  no admissible framework refutes $p$ on the shared base.
\end{definition}

Framework variance
(\DefRef{se100.def.FrameworkVariant}{Framework-Variant Proposition})
and framework invariance are stated over the same shared base
$\Substrate \cup \Framework$.
They are opposed in one direction:
if $p$ is framework-variant with respect to $\Substrate$ and $\Frameworks$,
then some admissible $\Framework$ gives
$\Substrate \cup \Framework \entails \neg p$.
Adding $p$ to that same base yields
$\Substrate \cup \Framework \cup \{p\} \entails \bot$,
so $p$ is not framework-invariant.

\begin{definition}[Framework-Compatible Commitment Set]
  \label{se100.def.FrameworkCompatibleCommitmentSet}
  A set of substrate-layer commitments $C$ is \emph{framework-compatible}
  with respect to $\Frameworks$ if, for every admissible framework
  $\Framework \in \Frameworks$,
  \[
    C \cup \Framework \notentails \bot .
  \]
\end{definition}

Framework invariance is a positive guarantee,
rather than absence of observed disagreement.
A proposition may fail to be framework-invariant
even when it is not framework-variant.
For example,
suppose no admissible framework derives $p$ on the shared base,
but some admissible framework refutes $p$ on that base:
\[
  \Substrate \cup \Framework' \entails \neg p .
\]
Then $p$ is not framework-variant, because no admissible framework derives $p$.
It is also not framework-invariant, because adding $p$ to the refuting base gives
\[
  \Substrate \cup \Framework' \cup \{p\} \entails \bot .
\]

The proposition-level definition tests whether a candidate proposition
can be added to the substrate while preserving compatibility with every
admissible framework.
The commitment-set definition tests whether a set of commitments is
compatible with every admissible framework.
Applied to $\SubstrateRef$,
the commitment-set definition gives the compatibility asserted by
\AssumpRef{se100.assump.ReferentialCommonGround}{Referential Common Ground}.

The neutrality constraint below requires this positive compatibility guarantee:
compatibility with every admissible framework,
including admissible frameworks not known when the substrate is designed.


% =============================================================================
\subsection*{Neutrality Conditions}
% =============================================================================


% ================  Contest ========================

\begin{definition}[Contested Causal or Normative Proposition]
  \label{se100.def.ContestedCausalNormative}
  A causal or normative proposition $p$ is \emph{contested}
  in the relevant accountability context
  if its acceptance, rejection, interpretation, or application
  is not fixed by the referential commitments $\SubstrateRef$,
  so that it may vary across admissible frameworks.

  Let $C_{cn}$ denote the class of contested causal or normative propositions.
\end{definition}

\begin{assumption}[Contestability]
  \label{se100.assump.Contestability}
  No $p \in C_{cn}$ is guaranteed framework-invariant at design time.
\end{assumption}

Contestability does not mean that propositions in $C_{cn}$ are
false, arbitrary, unknowable, or unsupported.
It means that the substrate cannot certify their universal compatibility
across admissible frameworks at design time.
Claims that a causal or normative proposition is
well supported,
evidence-backed,
legally required,
institutionally authorized,
or supported by a cited source
must themselves be represented as attributed, provenance-bearing assertions.
They are not exceptions to the restriction on substrate-layer commitment.

% ================  Common Ground ========================

\begin{assumption}[Referential Common Ground]
  \label{se100.assump.ReferentialCommonGround}
  Let $\SubstrateRef$ be the referential commitments of $\Substrate$
  (\DefRef{se100.def.ReferentialCommitments}{Referential Commitments}).
  For every admissible framework $\Framework$,
  \[
    \SubstrateRef \cup \Framework \notentails \bot
  \]
  and this holds for every admissible framework,
  including admissible frameworks not known when $\Substrate$ is designed.
\end{assumption}

\AssumpRef{se100.assump.ReferentialCommonGround}{Referential Common Ground}
states that $\SubstrateRef$ is a framework-compatible commitment set.
It formalizes the scope condition on which neutrality depends:
the frameworks in play contest interpretation, not reference.
The assumption requires the substrate's referential commitments to be compatible
with every admissible framework as a set, not just one by one.

This is an assumption, not a general property of records.
It can fail.
\SecRef{sec:boundary} treats its failure as the condition under which
neutrality is unavailable.

% permitted attribution is a structural restriction on the foundational layer,
% Referential common ground handles the common ground requirement
\begin{remark}[Attribution and Common Ground]
  \label{se100.remark.AttributionCommonGround}
  A permitted attribution proposition
  $\Asserts(x,\varphi)$ commits the substrate to the attribution by $x$,
  not to the asserted proposition $\varphi$
  (\DefRef{se100.def.PermittedAttributionProposition}{Permitted Attribution Proposition}).
  An admissible framework may reject $\varphi$ without rejecting that
  $x$ asserted it.

  A permitted attribution proposition need not be entailed by $\SubstrateRef$.
  It is a permitted foundational-layer commitment
  because its attributional basis is fixed by $\SubstrateRef$.
  Under
  \AssumpRef{se100.assump.ReferentialCommonGround}{Referential Common Ground},
  attributional bases fixed by $\SubstrateRef$ are common ground among
  admissible frameworks.
  Therefore, adding the permitted attribution proposition does not introduce
  a contested attributional basis into the foundational layer:
  for every admissible framework $\Framework$ compatible with
  $\SubstrateRef$,
  \[
    \SubstrateRef \cup \{\Asserts(x,\varphi)\} \cup \Framework
    \notentails \bot .
  \]

  If an admissible framework rejects $\varphi$,
  that does not by itself contest the attribution proposition $\Asserts(x,\varphi)$.
  But if admissible frameworks contest whether the attributional basis is
  fixed by $\SubstrateRef$, for example, who asserted the claim, what was
  asserted, or what provenance identifies the assertion occurrence,
  then the attribution proposition is not permitted at the foundational layer.
  That is a boundary case, not an exception to the neutrality constraint.
\end{remark}


\begin{definition}[Interpretive Non-Commitment]
  \label{se100.def.InterpretiveNonCommitment}
  A substrate $\Substrate$ satisfies \emph{interpretive non-commitment} if it
  makes no substrate-layer commitment to any proposition that is
  framework-variant with respect to $\Substrate$ and $\Frameworks$.
  Equivalently, if
  $p$ is framework-variant with respect to $\Substrate$ and $\Frameworks$,
  then
  $\Substrate \notentails p$ and $\Substrate \notentails \neg p$.
\end{definition}

Interpretive non-commitment does not prevent the substrate from recording that
a framework, source, agent, institution, record, or document asserts $p$.
It prevents the substrate from asserting $p$ independently of attribution.

\begin{definition}[Extension Stability]
  \label{se100.def.ExtensionStability}
  A substrate $\Substrate$ satisfies \emph{extension stability} if, for every
  admissible framework $\Framework$,
  \[
    \Substrate \cup \Framework \notentails \bot
  \]
\end{definition}

This is pairwise compatibility between the substrate and each admissible
framework, not global consistency across all admissible frameworks.
The distinction matters: admissible frameworks may contradict one another.
A neutral substrate need not make all frameworks mutually consistent.
It must remain separately usable by each admissible framework without causing contradiction.

% =========== Substrate Consistency ===========

\begin{assumption}[Substrate Consistency]
  \label{se100.assump.SubstrateConsistency}
  Every substrate considered in this paper is internally consistent:
  $\Substrate \notentails \bot$.
\end{assumption}

This assumption excludes degenerate substrates from the domain of analysis.
An internally inconsistent substrate cannot serve as the stable shared base
presupposed by the definitions.
The explicit assumption is retained to make the baseline condition independent
of whether any admissible framework has yet entered consideration.
It prevents vacuous cases in which neutrality conditions would hold only
because there are no admissible frameworks against which the substrate is tested.

% ======================   Remark: Relation Between the Two Properties ======================

\begin{remark}[Relation Between the Two Properties]
  \label{se100.remark.PropertyRelation}
  Extension stability entails interpretive non-commitment.
  To see this, suppose that $p$ is framework-variant with respect to
  $\Substrate$ and $\Frameworks$.
  Then there are admissible frameworks $\FrameworkOne, \FrameworkTwo \in \Frameworks$
  such that
  \[
    \Substrate \cup \FrameworkOne \entails p
    \qquad\text{and}\qquad
    \Substrate \cup \FrameworkTwo \entails \neg p .
  \]

  If $\Substrate \entails p$, then
  $\Substrate \cup \FrameworkTwo \entails p$
  and
  $\Substrate \cup \FrameworkTwo \entails \neg p$,
  so
  \[
    \Substrate \cup \FrameworkTwo \entails \bot .
  \]
  If $\Substrate \entails \neg p$, then
  $\Substrate \cup \FrameworkOne \entails \neg p$
  and
  $\Substrate \cup \FrameworkOne \entails p$,
  so
  \[
    \Substrate \cup \FrameworkOne \entails \bot .
  \]
  In either case, extension stability fails.
  Therefore, if extension stability holds, the substrate cannot commit to
  either side of a framework-variant proposition.

  This paper uses only this direction.
  The two properties are retained under separate names because they identify
  distinct design failures:
  substrate-layer commitment on disputed content and breakdown of layerability.
\end{remark}

% =============================================================================
\section{The Neutrality Constraint}
\label{sec:neutrality}
% =============================================================================

\begin{definition}[Neutrality by Design]
  \label{se100.def.NeutralityByDesign}
  A substrate $\Substrate$ is \emph{neutral}
  if it satisfies interpretive non-commitment
  (\DefRef{se100.def.InterpretiveNonCommitment}{Interpretive Non-Commitment})
  and extension stability
  (\DefRef{se100.def.ExtensionStability}{Extension Stability}).

  $\Substrate$ is \emph{neutral by design} if its neutrality is guaranteed
  at design time for every admissible framework in $\Frameworks$,
  including admissible frameworks not known when $\Substrate$ is designed
  (\NoteRef{se100.note.FrameworkClass}{The Framework Class $\Frameworks$}).
  The design-time guarantee must follow from
  the membership of the foundational layer in the permitted classes:
  the referential commitments $\SubstrateRef$ and permitted attribution propositions.
  It must not rest on enumerating admissible frameworks or
  on adopting any causal or normative proposition
  as an object-level substrate-layer commitment.
\end{definition}

The distinction is between a semantic property and a design-time guarantee.
Neutrality is the semantic property that interpretive non-commitment and
extension stability hold relative to $\Frameworks$.
Neutrality by design is stronger:
it requires those properties to be guaranteed by the structure of the substrate,
not by checking admissible frameworks one by one.

\begin{constraint}[Neutrality]
  \label{se100.constraint.Neutrality}
  Let $\Substrate$ be a substrate intended
  to remain usable by every admissible framework,
  including admissible frameworks not known when $\Substrate$ is designed.
  Assume
  \AssumpRef{se100.assump.Contestability}{Contestability} and
  \AssumpRef{se100.assump.ReferentialCommonGround}{Referential Common Ground}.
  Then $\Substrate$'s neutrality is guaranteed at design time
  in the sense of \DefRef{se100.def.NeutralityByDesign}{Neutrality by Design}
  if and only if its foundational layer is restricted to
  the referential commitments $\SubstrateRef$
  and permitted attribution propositions
  (\DefRef{se100.def.PermittedAttributionProposition}{Permitted Attribution Proposition}).
  In particular, $\Substrate$ makes no object-level causal or normative commitment
  (\DefRef{se100.def.ObjectLevelCausalNormativeCommitment}
  {Object-Level Causal or Normative Commitment}).
\end{constraint}

Informally:
if a shared record must remain usable by admissible frameworks that may
permanently disagree about causation and norms,
then the shared substrate may commit only
to referential commitments and permitted attribution propositions.
At the foundational layer,
causal and normative content may appear only through
permitted attribution propositions:
the substrate commits to the attribution,
not to the asserted causal or normative proposition.

Neutrality and neutrality by design are as in
\DefRef{se100.def.NeutralityByDesign}{Neutrality by Design}.
The constraint says that, under the two assumptions,
the restriction to referential commitments and permitted attribution propositions
is necessary and sufficient for neutrality by design.

This is a design constraint.
It is stated in full so that its dependence on the two assumptions is clear.

\emph{Sufficiency.}
Suppose the foundational layer of $\Substrate$ is restricted to
$\SubstrateRef$ and permitted attribution propositions.
By
\AssumpRef{se100.assump.ReferentialCommonGround}{Referential Common Ground},
$\SubstrateRef \cup \Framework \notentails \bot$
for every admissible $\Framework$.
By
\RemRef{se100.remark.AttributionCommonGround}{Attribution and Common Ground},
permitted attribution propositions do not require any admissible framework to
accept the asserted propositions and do not introduce contested attributional
basis into the foundational layer.
Since the foundational layer contains only
referential commitments and permitted attribution propositions,
$\Substrate \cup \Framework \notentails \bot$
for every admissible $\Framework$,
so extension stability holds.

By
\RemRef{se100.remark.PropertyRelation}{Relation Between the Two Properties},
extension stability entails interpretive non-commitment.
Thus $\Substrate$ is \emph{neutral}.

Because the guarantee follows from the structure of the foundational layer,
rather than from checking admissible frameworks one by one,
the substrate is \emph{neutral by design}.

Sufficiency uses the common-ground assumption essentially:
referential commitments and attributional bases are not inherently uncontestable.
Where that assumption fails, this direction fails
(\SecRef{sec:boundary}).

\emph{Design-time necessity.}
Suppose $\Substrate$ makes a substrate-layer commitment to a proposition $p$
outside the permitted classes:
$p$ is neither a referential commitment nor an attribution proposition
whose attributional basis is fixed by the substrate's referential commitments.
Referential commitments are exactly $\SubstrateRef$
(\DefRef{se100.def.ReferentialCommitments}{Referential Commitments}),
and attribution propositions are propositions of the form
$\Asserts(x,\varphi)$, for some identified source $x$
(\DefRef{se100.def.AttributionProposition}{Attribution Proposition}).

Since $p$ lies outside the permitted classes,
$p$ is not a referential commitment;
otherwise it would fall within $\SubstrateRef$.
There are two cases.

\emph{Case 1: $p \in C_{cn}$ and $p$ is not an attribution proposition.}
By
\AssumpRef{se100.assump.Contestability}{Contestability},
$p$ is not guaranteed framework-invariant at design time.
Thus the assumptions do not certify its compatibility
with every admissible framework.

\emph{Case 2: $p$ lies outside $C_{cn}$ or is an attribution proposition
  whose attributional basis is not fixed by $\SubstrateRef$.}
If $p$ is an attribution proposition
whose attributional basis is not fixed by $\SubstrateRef$, then
\AssumpRef{se100.assump.ReferentialCommonGround}{Referential Common Ground}
does not certify its compatibility with every admissible framework.
If the attributional basis itself is contested,
the boundary condition applies (\SecRef{sec:boundary}).

If $p$ is not an attribution proposition and not a referential commitment,
then its compatibility with every admissible framework is supplied by neither
\AssumpRef{se100.assump.ReferentialCommonGround}{Referential Common Ground}
nor \AssumpRef{se100.assump.Contestability}{Contestability}.

In both cases, the two assumptions certify compatibility only for
$\SubstrateRef$ and permitted attribution propositions.
For $p$ outside those classes,
neutrality does not follow from membership of
the foundational layer in the permitted classes.
Thus $\Substrate$ is not neutral by design in the sense of
\DefRef{se100.def.NeutralityByDesign}{Neutrality by Design}.

The representational mechanism that implements this layer separation is
reification of interpretive claims.

% =============================================================================
\section{Mechanism: Reification of Interpretive Claims}
\label{sec:mechanism}
% =============================================================================

The mechanism of neutrality is reification.
A non-neutral substrate asserts a
contested causal or normative proposition:
\[
  \Substrate \entails \varphi
\]
A neutral substrate instead represents
that some framework, source, agent, institution, record, or document $x$
asserts the proposition:
\[
  \Substrate \entails \Asserts(x,\varphi)
\]
The substrate commits to the attribution proposition
$\Asserts(x,\varphi)$,
not to the asserted proposition $\varphi$
(\DefRef{se100.def.SubstrateCommitment}{Substrate-Layer Commitment}).

This allows the record to preserve accountability-relevant content without
making the shared layer partisan among admissible frameworks.
The substrate may record:
the decision that was made;
the person or entity affected;
the time of the decision;
the instrument or system involved;
the policy or rule cited;
the source that asserted the claim;
the text of the claim;
the causal or normative proposition asserted;
and the provenance of the assertion.
The substrate must not convert the asserted proposition into a
substrate-layer commitment to that proposition.

% =============================================================================
\section{Worked Example}
\label{sec:example}
% =============================================================================

\begin{example}[Reification Fragment]
  \label{se100.example.ReificationFragment}
  Consider a decision record concerning an automated eligibility denial.
  The record involves a subject $u$, a decision $d$, a model or instrument $m$,
  a timestamp $t$, an institutional policy $r$, and
  a claim asserted by source framework $\Framework$.
\end{example}

The example presupposes referential common ground:
the frameworks in play share enough reference to identify
which subject, which decision, which instrument, and which policy the record
concerns.
What they contest is interpretation.

\textbf{Non-neutral form.}
The substrate makes object-level causal or normative commitments
(\DefRef{se100.def.ObjectLevelCausalNormativeCommitment}
{Object-Level Causal or Normative Commitment}):
\[
  \Substrate \entails \pred{caused\_by}(\pred{score}(u,m), \pred{denial}(d))
  \qquad
  \Substrate \entails \pred{justified\_by}(\pred{denial}(d), \pred{policy}(r))
\]
If admissible frameworks disagree about whether the score caused the denial,
whether the policy applied, or whether the denial was justified,
these commitments make the substrate non-neutral.

\textbf{Neutral form.}
The substrate makes referential commitments:
\begin{align*}
  \Substrate & \entails \pred{Subject}(u)
             & \Substrate                    & \entails \pred{occurred\_at}(d,t)     \\
  \Substrate & \entails \pred{Decision}(d)
             & \Substrate                    & \entails \pred{involved}(d,u)         \\
  \Substrate & \entails \pred{Instrument}(m)
             & \Substrate                    & \entails \pred{used\_instrument}(d,m) \\
  \Substrate & \entails \pred{Policy}(r)
             & \Substrate                    & \entails \pred{cites}(d,r)
\end{align*}
and records the contested content as attribution propositions:
\[
  \Substrate \entails \Asserts(\Framework,\;
  \pred{caused\_by}(\pred{score}(u,m), \pred{denial}(d)))
\]
\[
  \Substrate \entails \Asserts(\Framework,\;
  \pred{justified\_by}(\pred{denial}(d), \pred{policy}(r)))
\]
The substrate preserves that source framework $\Framework$
made these causal and normative claims.
It commits to the attribution propositions,
not to the asserted propositions.

\textbf{Predicate-level reading.}
The neutral record contains the
decision, subject, timestamp, identifiers, institutional artifacts, and
provenance-bearing attribution propositions.
It contains no unattributed top-level content whose meaning is equivalent to
the non-neutral commitments above.
This holds regardless of field naming:
a field named \texttt{reason}, \texttt{cause},
\texttt{basis}, \texttt{risk}, \texttt{fault}, or \texttt{violation}
is not prohibited by name.
The question is whether, in the schema and use context,
the field functions as a substrate-layer commitment to an object-level causal or normative proposition.
For example, a cited policy identifier or recorded model score may be referential content,
while an unattributed assertion that the policy justified the decision,
the score caused the denial, or the subject violated a rule is
causal or normative content on the substrate's own account.
Causal or normative vocabulary may appear;
the test is whether the substrate makes a substrate-layer commitment
to the asserted causal or normative proposition itself,
rather than to an attribution proposition about that proposition
(\NoteRef{se100.note.CausalNormative}{Causal and Normative Content as Primitive
  Classifications}).

The neutral form supports multiple downstream interpretations without
requiring the substrate to choose among them.
An institutional framework may treat the attributed claim as sufficient
for its purpose.
An audit framework may question whether the cited policy applied.
A claimant-side framework may contest the causal proposition.
A regulatory framework may evaluate whether the attributed claim satisfies
a disclosure rule.
The substrate remains usable by all of them
because it preserves the shared referential base and the attribution
propositions without making an object-level substrate-layer commitment
to the contested interpretation.

% =============================================================================
\section{Boundary Condition: When Neutrality Is Unavailable}
\label{sec:boundary}
% =============================================================================

Neutrality depends on
\AssumpRef{se100.assump.ReferentialCommonGround}{Referential Common Ground}.
That assumption can fail in two ways relevant here:
admissible frameworks may contest the substrate's referential regime,
or they may contest the attributional basis of an attribution proposition.

First, admissible frameworks may disagree about the substrate's
individuation, co-reference, or persistence conditions.
In that case, some admissible framework is incompatible with
$\SubstrateRef$.
\AssumpRef{se100.assump.ReferentialCommonGround}{Referential Common Ground}
fails, and neutrality is unavailable at the original foundational layer.

For example, one framework may treat two records as referring to the same
institutional event while another admissible framework treats them as distinct events.
One framework may treat an institutional identity as persistent across a merger
while another treats the merger as identity-breaking.
In such cases, the disagreement concerns the referential regime itself.
The substrate cannot be neutral among frameworks that do not share enough
reference to identify what is being interpreted.

Second, admissible frameworks may contest
the attributional basis of an attribution proposition.
A framework may reject the asserted proposition $\varphi$
while still accepting that source $x$ asserted $\varphi$.
That disagreement does not defeat neutrality.
But if admissible frameworks disagree about who made the assertion,
which assertion occurrence is being recorded,
what content was asserted,
or what provenance fixes the attribution,
then the attributional basis itself is contested.
In that case, the attribution proposition is not permitted at the
foundational layer.

These boundary cases do not make accountability impossible.
They mean the
contested referential regime or contested attributional basis
must be represented as an attributed claim,
or moved to a higher-level structure capable of tracking the disagreement.
What the original substrate cannot do is
treat one referential regime or one attributional basis
as neutral while admissible frameworks contest
the conditions on which that treatment depends.

This boundary condition prevents the neutrality constraint from applying
where reference or attributional basis is contested.
Without this condition,
a contested referential scheme or contested attributional basis
could be placed into the shared base as if it were neutral.
A neutral substrate must not let reference or attributional basis
carry an unattributed causal or normative commitment
(\DefRef{se100.def.ObjectLevelCausalNormativeCommitment}
{Object-Level Causal or Normative Commitment}).

% =============================================================================
\section{Scope and Limits}
\label{sec:non-goals}
% =============================================================================

\textbf{This paper does not propose a data format.}
A neutral substrate may be implemented using relational databases,
graph structures, event logs, RDF-like representations,
document stores, schemas, or application-specific formats.
The neutrality constraint applies to the substrate-layer commitments
made by the foundational layer,
not to any serialization or storage technology.

\textbf{This paper does not adjudicate among frameworks.}
Admissibility identifies minimum conditions for participation in the shared
record:
internal consistency,
evidentiary grounding,
and a documented interpretive function.
It does not decide which admissible framework should prevail in a dispute.

\textbf{This paper does not claim a neutral admissibility gate.}
As stated in \SecRef{sec:definitions},
the admissibility gate embeds design choices about what counts as
a framework that can enter the shared record.
The neutrality defined here is conditional on those choices.
The gate's role is to make participating frameworks identifiable,
scoped, and citable,
not to determine which framework is correct.

\textbf{This paper does not solve accountability generally.}
Neutrality is a structural condition for
shared records under persistent disagreement.
It does not determine liability, responsibility, institutional authority,
causal truth, legal compliance, policy adequacy, or
which interpretation should govern.
It makes such disputes representable without
forcing the shared substrate to decide them.

\textbf{This paper does not prohibit causal or normative content.}
Under the contestability assumption,
it prohibits substrate-layer commitments to
object-level causal or normative propositions
in the foundational layer.
Causal and normative claims remain representable by attribution with provenance.

% =============================================================================
\section{Discussion}
\label{sec:discussion}
% =============================================================================

The constraint is most useful in applied ontology settings
where reference is shared but interpretation remains persistently contested.
Examples include
litigation,
forensic discovery,
audit,
regulatory review,
institutional accountability records, and
multi-party AI-governance records.

Neutrality by design requires a specific representational discipline:
the substrate may assert reference,
but object-level causal or normative propositions may appear only through
attribution propositions.
This discipline matters in systems that combine records from multiple
institutions, jurisdictions, tools, or governance contexts,
where causal and normative terms often appear as ordinary metadata.
Neutral design does not prohibit such terms.
It asks what the substrate is committing to
(\DefRef{se100.def.SubstrateCommitment}{Substrate-Layer Commitment}).
If the substrate asserts that a denial was caused by a score,
that a policy justified a decision,
or that a party violated a norm,
the substrate has made an object-level causal or normative commitment.
If it records that a specified source made that claim under documented
provenance,
the substrate has preserved the claim without adopting it.

The constraint has practical value only where referential common ground holds.
In these settings, parties share enough reference
to identify which record,
which decision, which subject, and which instrument are at issue while disputing,
often persistently, what caused the outcome and whether it was justified.
This is a common condition of records under dispute.
The hard accountability cases in which reference itself is contested,
including disputed event identity and contested institutional persistence,
are the boundary cases of \SecRef{sec:boundary}.
This paper claims applicability where referential common ground holds
and unavailability where it does not.
It does not assume reference is always shared.

% =============================================================================
\section{Conclusion}
\label{sec:conclusion}
% =============================================================================

A shared accountability record must often serve parties who disagree
persistently about causation, norms, responsibility, legitimacy, or
institutional interpretation.
Such a record cannot remain neutral if its foundational layer asserts one
contested interpretation as an object-level substrate-layer commitment.

This does not make the record empty.
It makes the record usable for accountability.
The substrate can preserve who said what,
about which entity, event, decision, instrument, policy, or artifact,
under what provenance,
and within what scope.
It must not convert an attributed interpretation into
an object-level causal or normative commitment of the substrate.

Subsequent work will refine the referential regimes that determine when
referential common ground holds and
will specify checks for verifying the constraint in concrete record systems.

Neutrality is unavailable when the referential regime itself is contested.
But where enough reference can be shared,
a neutral substrate provides a structural base for disagreement without erasure,
accountability without premature adjudication,
and interpretation without unattributed commitment.

% =============================================================================
\section*{Statements and Declarations}
% =============================================================================

\subsection*{Author Contributions}
The author is the sole contributor to this work and
is responsible for all aspects of the research, authorship, and publication.

\subsection*{Use of AI-Assisted Tools}
AI-assisted tools were used for editing, formatting, and consistency checking.
The author reviewed all suggestions and is solely responsible for the content.

\subsection*{Declaration of Conflicting Interest}
The author declares no potential conflicts of interest with respect to the
research, authorship, and publication of this work.

\begin{thebibliography}{99}

  \bibitem[Carroll et~al.(2005)]{carroll2005named}
  Carroll, J.~J., Bizer, C., Hayes, P., and Stickler, P. (2005).
  Named graphs, provenance and trust.
  In \emph{Proceedings of the 14th International Conference on World Wide Web
    (WWW 2005)}, pages 613--622. ACM.

  \bibitem[de~Kleer(1986)]{dekleer1986atms}
  de~Kleer, J. (1986).
  An assumption-based TMS.
  \emph{Artificial Intelligence}, 28(2):127--162.

  \bibitem[Dung(1995)]{dung1995}
  Dung, P.~M. (1995).
  On the acceptability of arguments and its fundamental role in nonmonotonic
  reasoning, logic programming and $n$-person games.
  \emph{Artificial Intelligence}, 77(2):321--357.

  \bibitem[Floridi(2008)]{floridi2008loa}
  Floridi, L. (2008).
  The method of levels of abstraction.
  \emph{Minds and Machines}, 18(3), 303--329.
  doi:10.1007/s11023-008-9113-7.

  \bibitem[Gangemi and Mika(2003)]{gangemi2003descriptions}
  Gangemi, A., and Mika, P. (2003).
  Understanding the Semantic Web through descriptions and situations.
  In R. Meersman, Z. Tari, and D. C. Schmidt (Eds.),
  \emph{On the Move to Meaningful Internet Systems 2003:
    CoopIS, DOA, and ODBASE},
  Lecture Notes in Computer Science, vol. 2888.
  Springer.
  doi:10.1007/978-3-540-39964-3\_44.

  \bibitem[Guarino(1998)]{guarino1998}
  Guarino, N. (1998).
  Formal ontology and information systems.
  In \emph{Formal Ontology in Information Systems (FOIS~'98)}, pages 3--15.
  IOS Press.

  \bibitem[Hayes(2004)]{hayes2004rdf}
  Hayes, P. (2004).
  \emph{RDF Semantics}.
  W3C Recommendation, World Wide Web Consortium.

  \bibitem[Lebo et~al.(2013)]{lebo2013prov}
  Lebo, T., Sahoo, S., and McGuinness, D.~L., editors (2013).
  \emph{PROV-O: The PROV Ontology}.
  W3C Recommendation, World Wide Web Consortium.

  \bibitem[McCarthy(1993)]{mccarthy1993}
  McCarthy, J. (1993).
  Notes on formalizing context.
  In \emph{Proceedings of the 13th International Joint Conference on Artificial
    Intelligence (IJCAI-93)}, pages 555--560.

  \bibitem[Searle(1969)]{searle1969}
  Searle, J.~R. (1969).
  \emph{Speech Acts: An Essay in the Philosophy of Language}.
  Cambridge University Press.

\end{thebibliography}

\end{document}
